\section{公司的解散、清算}
\begin{enumerate}
  \item 一般解散。
  \item 股东大会决议解散。
  \item 合并分立。
  \item 强制解散。
  \item 股东请求法院解散。公司僵局。
        \begin{enumerate}
          \item 经营管理发生严重困难。表决权 10\% 以上股东。
          \item 2 年无法开股东会；能开会但不能做决议。人的矛盾，治理困境。
          \item 解散和清算不并行，相分离。
          \item 注重调解。
        \end{enumerate}
  \item 清算。终结已解散公司的一切法律关系，处理公司剩余财产的程序。
        \begin{enumerate}
          \item 自行清算。
          \item 强制清算。法院主导。
                \begin{enumerate}
                  \item 3 种情形。逾期、拖延、违法。
                  \item 4 种主体。债权人、股东、董事、其他厉害关系人。
                \end{enumerate}
          \item 清算组。
          \item 清算向破产转换。资不抵债。
          \item 法律责任。
                \begin{enumerate}
                  \item 未成立清算组。逾期造成损失，在损失范围内承担责任。
                  \item 怠于履行义务，毁损灭失，无法清算，连带责任。抗辩：采取了积极措施。
                \end{enumerate}
        \end{enumerate}
\end{enumerate}